\section{The three 7/8 challenge statements}\label{app:78-listings}

Listing~\ref{lst:78-challenges} quotes the statements of the QuasiRiemannHypothesis, DirichletSevenEighths and HeckeSevenEighths challenges at openai/math fd4aeeb.

\begin{listing}[tbp]
\caption{The 7/8 challenges at openai/math fd4aeeb: lean/ComparatorChallenges/QuasiRiemannHypothesis.lean:5-6 (sha256 065f8c9a\ldots), DirichletSevenEighths.lean:8-11 (sha256 8fb13ad9\ldots) and HeckeSevenEighths.lean:237-239 (sha256 cb5e404f\ldots).}\label{lst:78-challenges}
\begin{Verbatim}
theorem riemannZeta_ne_zero_of_seven_eighths_lt_re
    {s : ℂ} (hs : (7 / 8 : ℝ) < s.re) : riemannZeta s ≠ 0 := by
\end{Verbatim}
\begin{Verbatim}
theorem LFunction_ne_zero_of_seven_eighths_lt_re
    {q : ℕ} [NeZero q] (χ : DirichletCharacter ℂ q) {s : ℂ}
    (hs : (7 / 8 : ℝ) < s.re) (hpole : ¬ (χ = 1 ∧ s = 1)) :
    _root_.DirichletCharacter.LFunction χ s ≠ 0 := by
\end{Verbatim}
\begin{Verbatim}
theorem LFunction_ne_zero_of_seven_eighths_lt_re
    (χ : Character) {s : ℂ} (hs : (7 / 8 : ℝ) < s.re)
    (hpole : ¬ (χ.residue = 1 ∧ s = 1)) : LFunction χ s ≠ 0 := by sorry
\end{Verbatim}
\end{listing}
